% ---
% title: 洛陽伽藍記序
% author: 楊衒之
% ---
\begin{guji}[20]
洛陽伽藍記序

\begin{flushright}
魏撫軍府司馬楊衒之撰
\end{flushright}

三墳五典之說，\jiazhu{孔安國尚書序：伏犧、神農、黃帝之書，謂之三墳；少昊、顓頊、高辛、唐、虞之書，謂之五典。}
九流百代之言，\jiazhu{漢書藝文志列儒、道、陰陽、法、名、墨、縱橫、雜、農，是為九流。}
並理在人區，而義兼天外。
至於一乘二諦之原，\jiazhu{一乘，佛乘也。二諦，真諦、俗諦也。}
三明六通之旨，\jiazhu{三明，宿命、天眼、漏盡。六通，於三明之外，加天耳、他心、神足。}
西域備詳，東土靡記。
自頃日感夢，滿月流光，\jiazhu{謂漢明帝夜夢金人，頂有光明，遂遣使天竺問其道術。事見後漢書西域傳。}
陽門飾毫眉之像，夜臺圖紺髮之形。\jiazhu{牟子理惑論：於開陽城門上作佛像；明帝豫修壽陵，曰顯節，亦於其上作佛圖像。毫眉、紺髮，皆佛相也。}
邇來奔競，其風遂廣。
至晉永嘉，唯有寺四十二所。\jiazhu{永嘉，晉懷帝年號。}
逮皇魏受圖，光宅嵩洛，\jiazhu{謂高祖孝文帝太和十七年遷都洛陽。}
篤信彌繁，法教愈盛。
王侯貴臣，棄象馬如脫屣；\jiazhu{史記封禪書：吾視去妻子如脫屣耳。言棄之易也。}
庶士豪家，捨資財若遺跡。
於是昭提櫛比，\jiazhu{昭提，即招提，寺之別名。}
寶塔駢羅，爭寫天上之姿，競摹山中之影。
金刹與靈臺比高，\jiazhu{靈臺，漢光武所築，在洛陽城南。}
講殿共阿房等壯。\jiazhu{阿房，秦宮名。}
豈直木衣綈繡，土被朱紫而已哉！

暨永熙多難，皇輿遷鄴，\jiazhu{永熙，孝武帝年號。永熙三年，帝西奔長安；孝靜帝即位，遷都於鄴。}
諸寺僧尼，亦與時徙。
至武定五年，歲在丁卯，\jiazhu{武定，孝靜帝年號。五年當梁太清元年。}
余因行役，重覽洛陽。
城郭崩毀，宮室傾覆，寺觀灰燼，廟塔丘墟。
牆被蒿艾，巷羅荊棘。
野獸穴於荒階，山鳥巢於庭樹。
遊兒牧豎，躑躅於九逵；\jiazhu{爾雅釋宮：九達謂之逵。}
農夫耕老，藝黍於雙闕。\jiazhu{闕，宮門兩觀也。}
麥秀之感，非獨殷墟；\jiazhu{史記宋微子世家：箕子朝周，過故殷虛，感宮室毀壞，生禾黍，乃作麥秀之詩。}
黍離之悲，信哉周室！\jiazhu{毛詩序：黍離，閔宗周也。周大夫行役，至于宗周，過故宗廟宮室，盡為禾黍。}
京城表裏，凡有一千餘寺，今日寮廓，鐘聲罕聞。\jiazhu{寮廓，空曠貌。}
恐後世無傳，故撰斯記。
然寺數最多，不可遍寫，今之所錄，止大伽藍；\jiazhu{伽藍，梵言僧伽藍摩之略，華言眾園，僧所居也。}
其中小者，取其祥異，世諦俗事，因而出之。
先以城內為始，次及城外，表列門名，以記遠近，凡為五篇。\jiazhu{城內、城東、城南、城西、城北，各為一卷。}
余才非著述，多有遺漏，後之君子，詳其闕焉。

太和十七年，高祖遷都洛陽，詔司空公穆亮營造宮室。\jiazhu{穆亮，字幼輔，時為司空。}
洛陽城門依魏晉舊名。

東面有三門。北頭第一門曰建春門，漢曰上東門。阮籍詩曰「步出上東門」是也。\jiazhu{阮籍詠懷詩：步出上東門，北望首陽岑。}
魏晉曰建春門，高祖因而不改。
次南曰東陽門，漢曰中東門，魏晉曰東陽門，高祖因而不改。
次南曰青陽門，漢曰望京門，魏晉曰清明門，高祖改為青陽門。

南面有四門。東頭第一門曰開陽門。
初，漢光武遷都洛陽，作此門始成，而未有名，忽夜中有柱自來在樓上。
後琅邪郡開陽縣言南門一柱飛去，使來視之，則是也。\jiazhu{開陽，漢縣，屬琅邪郡。}
遂以開陽為名。自魏及晉因而不改，高祖亦然。
次西曰平昌門，漢曰平門，魏晉曰平昌門，高祖因而不改。
次西曰宣陽門，漢曰小苑門，魏晉曰宣陽門，高祖因而不改。
次西曰津陽門，漢曰津門，魏晉曰津陽門，高祖因而不改。

西面有四門。南頭第一門曰西明門，漢曰廣陽門，魏晉因而不改，高祖改為西明門。
次北曰西陽門，漢曰雍門，魏晉曰西明門，高祖改為西陽門。
次北曰閶闔門，漢曰上西門，上有銅璇璣玉衡，以齊七政。\jiazhu{尚書舜典：在璿璣玉衡，以齊七政。七政，日月五星也。}
魏晉曰閶闔門，高祖因而不改。
次北曰承明門。承明者，高祖所立，當金墉城前東西大道。\jiazhu{金墉城，魏明帝所築，在洛陽城西北隅。}
遷京之始，宮闕未就，高祖住在金墉城。
城西有王南寺，高祖數詣寺沙門論議，故通此門，而未有名，世人謂之新門。
時王公卿士常迎駕於新門。
高祖謂御史中尉李彪曰：\jiazhu{李彪，字道固。}
「曹植詩云：『謁帝承明廬。』此門宜以承明為稱。」\jiazhu{曹植贈白馬王彪詩：謁帝承明廬，逝將歸舊疆。}
遂名之。

北面有二門。西頭曰大夏門，漢曰夏門，魏晉曰大夏門。
嘗造三層樓，去地二十丈。
洛陽城門樓皆兩重，去地百尺，惟大夏門甍棟干雲。\jiazhu{甍，屋棟也。}
東頭曰廣莫門，漢曰穀門，魏晉曰廣莫門，高祖因而不改。
自廣莫門以西，至於大夏門，宮觀相連，被諸城上也。

一門有三道，所謂九軌。\jiazhu{周禮考工記：經涂九軌。}
\end{guji}

\section*{排印说明}
\begin{enumerate}
  \item 正文录《洛阳伽蓝记》卷首杨衒之自序全文，自“三坟五典之说”至“所谓九轨”：前半叙佛法东来与洛阳寺塔兴废，后半依东、南、西、北四面历叙洛阳十三城门。文字据通行本，诸本字句互有出入。
  \item 版式仿卷子本：竖排，自右而左，每行二十字，行间施乌丝栏，上下各有栏线。页面宽于屏幕时可左右拖动，从右端读起。
  \item 注为双行小字夹注，随文附于句下，文字为本页所拟，非旧注。注长则跨行接续，不留空白。
  \item 不用新式标点。句末于字旁施圈（句），句中施点（读）；引号、书名号一概不用，书名、篇名径书。
  \item 排版由 hatex 完成：\verb|\begin{guji}[20]| 竖排加栏，方括号内为每行字数；\verb|\jiazhu{…}| 为双行夹注；标点在 \texttt{guji} 内自动改为句读，源文件照常书写。
\end{enumerate}
